\documentclass[11pt]{article}
\usepackage[a4paper,margin=27mm]{geometry}
\usepackage{amsmath,amssymb,mathtools}
\usepackage{fontspec,unicode-math}
\setmainfont{Latin Modern Roman}
\setmathfont{Latin Modern Math}
\usepackage{xurl}
\usepackage[hidelinks]{hyperref}
\newcommand{\sref}[2]{\hyperlink{source:#1:#2}{[S#2]}}
\newcommand{\cataloguescope}[1]{\paragraph{Mathematical question.}#1}
\hypersetup{pdftitle={Problem statement}}
\begin{document}
\section{Erdős–Mollin–Walsh conjecture}\label{statement}

\subsection{Definitions and mathematical statement}

\paragraph{Shared notation.}$\mathbb N=\{1,2,\ldots\}$, and $\mathbb Z,\mathbb Q,\mathbb R,\mathbb C$ denote the integers, rationals, reals and complex numbers. Any other conventions are specified locally.


A positive integer $n$ is powerful when every prime divisor $p\mid n$ satisfies $p^2\mid n$. The integer $1$ is powerful by the empty condition. Equivalently, $n=a^2b^3$ for some positive integer $a$ and squarefree positive integer $b$.
\paragraph{Statement recorded in the supplied source.}
Is there no $N\in\mathbb N$ for which $N$, $N+1$ and $N+2$ are all powerful? \sref{statement}{2}
\subsection{Short English statement}
Can three consecutive positive integers ever all have every prime factor appearing to exponent at least two?
\subsection{Sources}
\begin{description}
\item[\hypertarget{source:statement:1}{S1}] ULAM.AI and the UnsolvedMath Contributors, UnsolvedMath: A Curated Collection of Open Mathematics Problems, catalogue entry. CC BY 4.0. \url{https://huggingface.co/datasets/ulamai/UnsolvedMath}.
\item[\hypertarget{source:statement:2}{S2}] Tsz Ho Chan, \emph{A note on three consecutive powerful numbers}, arXiv:2503.21485v1 (2025) Introduction, definition and Conjecture 1, p. 1. \url{https://arxiv.org/pdf/2503.21485}.
\item[\hypertarget{source:statement:3}{S3}] Wikipedia contributors, \emph{List of unsolved problems in mathematics}, revision 1366636767 (2026) Number theory, Erdős–Mollin–Walsh conjecture. \url{https://en.wikipedia.org/w/index.php?title=List_of_unsolved_problems_in_mathematics&oldid=1366636767#Number_theory}.
\end{description}
\label{end:statement}
\end{document}
